\documentclass{article}
\usepackage[UTF8]{ctex}
\usepackage{amsmath,amssymb,amsfonts}
\usepackage{enumitem}
\usepackage{geometry}
\geometry{a4paper, margin=1in}
% === 错题本排版优化（不参与题目标识） ===
\setlength{\parindent}{0pt}
\setlength{\parskip}{0.35em plus 0.10em minus 0.08em}
\setlength{\emergencystretch}{3em}
\setlist{itemsep=0.25em, topsep=0.40em, parsep=0pt}
\allowdisplaybreaks[2]
% === 错题本排版优化结束 ===
\begin{document}
\section*{中南大学2024年高等代数真题}
\section*{135．中南大学2024年高等代数试卷}

\subsection{题目1}
一．（16 分）证明：任意给定 \(k(k \geq 1)\) 个两两不等的素数 \(a_{1}, a_{2}, \cdots, a_{k}\) ，对任意的正整数 \(n(n>1)\) ，都有 \(\sqrt[n]{a_{1} a_{2} \cdots a_{k}}\) 是无理数。

\subsection{题目2}
二．（20 分）解答如下问题：
1．（10 分）设 \(n\) 阶行列式
\[
D=\left|\begin{array}{ccccc}
1 & -1 & -1 & \cdots & -1 \\
1 & 1 & -1 & \cdots & -1 \\
1 & 1 & 1 & \cdots & -1 \\
\vdots & \vdots & \vdots & & \vdots \\
1 & 1 & 1 & \cdots & 1
\end{array}\right|
\]
求 \(D\) 展开式的正项总数。
2．（10 分）计算 \(n\) 阶行列式
\[
D_{n}=\left|\begin{array}{cccc}
\frac{1-a_{1}^{n} b_{1}^{n}}{1-a_{1} b_{1}} & \frac{1-a_{1}^{n} b_{2}^{n}}{1-a_{1} b_{2}} & \cdots & \frac{1-a_{1}^{n} b_{n}^{n}}{1-a_{1} b_{n}} \\
\frac{1-a_{2}^{n} b_{1}^{n}}{1-a_{2} b_{1}} & \frac{1-a_{2}^{n} b_{2}^{n}}{1-a_{2} b_{2}} & \cdots & \frac{1-a_{2}^{n} b_{n}^{n}}{1-a_{2} b_{n}} \\
\vdots & \vdots & & \vdots \\
\frac{1-a_{n}^{n} b_{1}^{n}}{1-a_{n} b_{1}} & \frac{1-a_{n}^{n} b_{2}^{n}}{1-a_{n} b_{2}} & \cdots & \frac{1-a_{n}^{n} b_{n}^{n}}{1-a_{n} b_{n}}
\end{array}\right|
\]

\subsection{题目3}
三。（16 分）设 \(A\) 为 \(m \times n\) 实矩阵，\(b\) 是 \(m \times 1\) 实列向量。

（1） 证明： \(\operatorname{rank}\left(A^{\mathrm{T}} A\right)=\operatorname{rank}(A)\) ．
（2） 设 \(X=\left(x_{1}, x_{2}, \cdots, x_{n}\right)^{\mathrm{T}}\) ，证明：线性方程组 \(A^{\mathrm{T}} A X=A^{\mathrm{T}} b\) 有解。
（3） 试举反例说明对复矩阵 \(A\) ，结论 \(\operatorname{rank}\left(A^{\mathrm{T}} A\right)=\operatorname{rank}(A)\) 不一定成立，（2）中线性方程组不一定有解。

\subsection{题目4}
四．（16 分）设 \(W_{1}, W_{2}\) 是线性空间 \(P^{2 \times 2}\) 的子空间，\(W_{1}=L\left(A_{1}, A_{2}\right), W_{2}=L\left(B_{1}, B_{2}\right)\) ，其中
\[
A_{1}=\left(\begin{array}{ll}
1 & 1 \\
0 & 0
\end{array}\right), A_{2}=\left(\begin{array}{ll}
1 & 0 \\
1 & 1
\end{array}\right), B_{1}=\left(\begin{array}{ll}
0 & 0 \\
1 & 1
\end{array}\right), B_{2}=\left(\begin{array}{ll}
0 & 1 \\
1 & 0
\end{array}\right) .
\]

（1） 求 \(W_{1}+W_{2}\) 的一组基与维数。
（2） 求 \(W_{1} \cap W_{2}\) 的一组基与维数。
（3） 问 \(P^{2 \times 2}=W_{1} \oplus W_{2}\) 是否成立，为什么？

\subsection{题目5}
五。（18 分）给定前 \(n\) 个自然数 \(1,2, \cdots, n\) 的一个全排列 \(j_{1} j_{2} \cdots j_{n}\) ，定义复数域上线性空间 \(\mathbb{C}^{n \times n}\) 的一个线性变换 \(\varphi\) 如下：
\[
\varphi\left(\begin{array}{cccc}
a_{11} & a_{12} & \cdots & a_{1 n} \\
a_{21} & a_{22} & \cdots & a_{2 n} \\
\vdots & \vdots & & \vdots \\
a_{n 1} & a_{n 2} & \cdots & a_{n n}
\end{array}\right)=\left(\begin{array}{cccc}
a_{1 j_{1}} & a_{1 j_{2}} & \cdots & a_{1 j_{n}} \\
a_{2 j_{1}} & a_{2 j_{2}} & \cdots & a_{2 j_{n}} \\
\vdots & \vdots & & \vdots \\
a_{n j_{1}} & a_{n j_{2}} & \cdots & a_{n j_{n}}
\end{array}\right) .
\]

（1） 试确定 \(\varphi\) 的 \(n\) 个线性无关的特征向量。
（2） 证明：存在正整数 \(k\) ，使得对 \(\varphi\) 的任意特征值 \(\lambda\) ，都有 \(\lambda^{k}=1\) 。
（3） 证明：如果全排列 \(j_{1} j_{2} \cdots j_{n}=23 \cdots n 1\) ，那么相应的线性变换 \(\varphi\) 是可对角化的．

\subsection{题目6}
六。（16 分）设 \(A\) 为 \(n\) 阶半正定矩阵，\(E\) 为 \(n\) 阶单位矩阵。证明：
\[
|A+2024 E| \geq 2024^{n} .
\]
并且等号成立当且仅当 \(A=O\) ．

\subsection{题目7}
七．（16 分）设 \(\sigma, \tau\) 是 \(n\) 维线性空间 \(V\) 上的线性变换，\(\sigma^{2}=\sigma\) ．证明：

（1） \(V=\operatorname{Im} \sigma \oplus \operatorname{Ker} \sigma\) ．
（2） \(\sigma\) 的像 \(\operatorname{Im} \sigma\) 与核 \(\operatorname{Ker} \sigma\) 都是 \(\tau\) 的不变子空间当且仅当 \(\sigma \tau=\tau \sigma\) ．

\subsection{题目8}
八．（16 分）设数域 \(P\) 上的矩阵 \(A=\left(\begin{array}{ccc}1 & -3 & 3 \\ -2 & -6 & 13 \\ -1 & -4 & 8\end{array}\right)\) ．

（1） 求矩阵 \(A\) 的 Jordan标准形 \(J\) 。
（2） 求可逆矩阵 \(P\) ，使得 \(P^{-1} A P=J\) 。
（3） 求 \(A\) 的最小多项式．

\subsection{题目9}
九．（16 分）设 \(V\) 是实数域上所有 \(n\) 阶对称矩阵构成的线性空间，对任意的 \(A, B \in V\) ，定义
\[
(A, B)=\operatorname{tr}(A B) .
\]
其中 \(\operatorname{tr}(A B)\) 表示 \(A B\) 的迹．

（1） 证明：\(V\) 构成一欧氏空间。
（2） 求使 \(\operatorname{tr}(A)=0\) 的子空间 \(S\) 的维数．
（3） 求 \(S\) 的正交补空间 \(S^{\perp}\) 的维数。
\end{document}
